% =====================================================================
% TEMPORARY SKELETON - replace preamble/front matter with the official
% campus template once it is placed in thesis/template-kampus/.
% Chapter files (bab*.tex) are kept separate so they can be reused.
% =====================================================================
\documentclass[12pt,a4paper,oneside]{report}

\usepackage[utf8]{inputenc}
\usepackage[T1]{fontenc}
\usepackage[bahasa]{babel}
\usepackage[a4paper,left=4cm,right=3cm,top=3cm,bottom=3cm]{geometry} % [ISI NANTI: follow campus template]
\usepackage{setspace}
\usepackage{graphicx}
\usepackage{amsmath,amssymb}
\usepackage{booktabs}
\usepackage{csquotes}
\usepackage[backend=biber,style=authoryear,maxcitenames=2]{biblatex} % [ISI NANTI: citation style per campus/advisor]
\usepackage[hidelinks]{hyperref}

\addbibresource{refs.bib}
\graphicspath{{figures/}{../results/}}
\onehalfspacing

\title{KLASIFIKASI KONDISI SEL CD4 PADA INFEKSI HIV BERDASARKAN CITRA IMAGING FLOW CYTOMETRY MENGGUNAKAN MACHINE LEARNING}
\author{Rahul \\ NIM: [ISI NANTI]}
\date{[ISI NANTI]}

\begin{document}

\maketitle

\chapter*{INTISARI}
\addcontentsline{toc}{chapter}{INTISARI}
[ISI NANTI]

\chapter*{ABSTRACT}
\addcontentsline{toc}{chapter}{ABSTRACT}
[ISI NANTI]

\tableofcontents
\listoffigures
\listoftables

\chapter{PENDAHULUAN}\label{pendahuluan}
% Structure follows template-kampus/Bab1.tex.

\section{Latar Belakang}\label{latar-belakang}
% Gambaran permasalahan/kebutuhan dan alasan mengapa penelitian ini penting.
[ISI NANTI]

\section{Perumusan Masalah}\label{perumusan-masalah}
% (1) identifikasi akar permasalahan/kebutuhan, (2) pendekatan penyelesaiannya,
% termasuk pembatasan pendekatan yang dianggap layak dalam kerangka skripsi.
[ISI NANTI]

\section{Tujuan}\label{tujuan}
[ISI NANTI]

\section{Manfaat}\label{manfaat}
[ISI NANTI]

\chapter{TINJAUAN PUSTAKA DAN DASAR TEORI}

\section{Tinjauan Pustaka}
% TODO: penelitian terdahulu. Hanya sitasi dari refs.bib yang sudah diverifikasi (lihat docs/literatur/).
[ISI NANTI]

\section{Dasar Teori}

\subsection{Sel T CD4 dan Infeksi HIV}
[ISI NANTI]

\subsection{Imaging Flow Cytometry}
[ISI NANTI]

\subsection{Morfometri dan Fitur Bio-Optik Sel}
[ISI NANTI]

\subsection{Machine Learning untuk Klasifikasi Citra}
[ISI NANTI]

\subsection{Explainable AI (Grad-CAM)}
[ISI NANTI]

\subsection{Metrik Evaluasi}
[ISI NANTI]

\chapter{METODOLOGI PENELITIAN}

\section{Alat dan Bahan}
% TODO: perangkat keras/lunak, sumber dataset.
[ISI NANTI]

\section{Alur Penelitian}
[ISI NANTI]

\section{Dataset}
% TODO: sumber, jumlah citra per kelas (sehat/terinfeksi/lisis), kanal, cara pelabelan.
[ISI NANTI]

\section{Praproses Data}
[ISI NANTI]

\section{Ekstraksi Fitur}
[ISI NANTI]

\section{Model Klasifikasi}
[ISI NANTI]

\section{Skema Pelatihan dan Pembagian Data}
[ISI NANTI]

\section{Evaluasi dan Explainable AI}
[ISI NANTI]

\chapter{HASIL DAN PEMBAHASAN}\label{hasil-pembahasan}
% RULE: every number in this chapter must come from a file in results/.
% Guideline: interpret WHY the phenomena occur and compare with recent literature.

\section{Hasil}\label{hasil}
[ISI NANTI]

\section{Pembahasan}\label{pembahasan}
[ISI NANTI]

\chapter{PENUTUP}

\section{Kesimpulan}
[ISI NANTI]

\section{Saran}
[ISI NANTI]


\printbibliography[heading=bibintoc,title={DAFTAR PUSTAKA}]

\end{document}
